\chapter{Excoriation}

\section*{Definition and Overview}
Excoriation is the medical term for the skin damage caused by scratching, rubbing, or picking. It appears as scratches, raw areas, and scabs, and it is usually a symptom of an underlying condition that causes itching, such as eczema, insect bites, or dry skin. Excoriation can also be a mental health condition in itself, called excoriation (skin-picking) disorder or dermatillomania, in which a person repeatedly picks at their skin, causing tissue damage and distress. Treating the underlying cause of the itching is the key to healing excoriated skin.

\section*{Epidemiology}
Excoriated skin is extremely common, since many skin conditions cause itching that leads to scratching. Skin-picking disorder affects a smaller but significant proportion of people, with estimates of 1 to 2 per cent, and it is more common in women. Skin picking often begins in adolescence, and it can cause significant distress, scarring, and infection, yet many people do not seek help because of shame.

\section*{Aetiology and Pathophysiology}
Excoriation results from the normal response to itch, or from compulsive picking.
\begin{itemize}
    \item \textbf{Itchy skin conditions:} eczema, psoriasis, scabies, insect bites, and dry skin cause itching, and scratching damages the skin
    \item \textbf{The itch-scratch cycle:} scratching damages the skin, which releases more itch chemicals, leading to more scratching
    \item \textbf{Nerve and dry skin:} dry skin and nerve problems can cause itch without a visible rash
    \item \textbf{Skin-picking disorder:} a mental health condition in which picking is driven by tension, habit, or the urge to remove perceived flaws
    \item \textbf{Associated conditions:} skin picking commonly coexists with anxiety, depression, obsessive-compulsive disorder, and body-focused repetitive behaviours
    \item \textbf{Effects:} broken skin is vulnerable to infection and leaves scars and pigment changes
\end{itemize}

\section*{Clinical Features}
The appearance of excoriation depends on its cause.
\begin{itemize}
    \item Scratches, linear marks, and raw areas of skin
    \item Scabs and crusts over the damaged areas
    \item Itching, which may be intense and worse at night
    \item Dry, rough skin around the excoriations
    \item In skin-picking disorder: repetitive picking at the face, arms, or other areas, often in response to minor imperfections
    \item Difficulty stopping the behaviour, despite wanting to
    \item Scarring and changes in skin colour over time
\end{itemize}

\section*{Differential Diagnosis}
Other skin damage can resemble excoriation.
\begin{itemize}
    \item Impetigo and other bacterial infections, with crusting and honey-coloured scabs
    \item Herpes infections, with clusters of blisters
    \item Dermatitis artefacta, in which skin damage is caused deliberately and secretly
    \item Folliculitis and other inflammatory skin conditions
    \item Healing wounds from other causes
\end{itemize}

\section*{When to Seek Medical Care}
People with persistent itching and scratching, or with skin picking they cannot control, should seek help from a GP. Treatment of the underlying skin condition or mental health problem is effective.

\textbf{Call 000 (or local emergency services) for:}
\begin{itemize}
    \item Rapidly spreading redness, pain, and fever, which may indicate a serious skin infection
    \item Signs of sepsis, including confusion and severe illness
    \item Bleeding that does not stop
    \item A deep wound requiring medical attention
\end{itemize}

\section*{Diagnosis and Investigations}
Diagnosis identifies the cause of the excoriation.
\begin{itemize}
    \item \textbf{History:} the pattern of itching, scratching, and picking, and any triggers
    \item \textbf{Examination:} the distribution and appearance of the skin damage
    \item \textbf{Skin assessment:} identification of eczema, scabies, or other underlying conditions
    \item \textbf{Mental health assessment:} for skin-picking disorder, including the impact on daily life
    \item \textbf{Swabs:} if infection is suspected
\end{itemize}

\section*{Management}
Management treats the underlying cause and protects the skin.
\begin{itemize}
    \item \textbf{Treating itchy skin:} moisturisers, topical steroids, and antihistamines for eczema and other itchy conditions
    \item \textbf{Treating specific causes:} scabies treatment, insect bite relief, and management of dry skin
    \item \textbf{Breaking the itch-scratch cycle:} cooling, distraction, and short nails
    \item \textbf{Skin-picking disorder:} cognitive behavioural therapy, including habit reversal training
    \item \textbf{Medicines:} antidepressants may help some people with skin picking
    \item \textbf{Wound care:} keeping the skin clean, and treating any infection
    \item \textbf{Scar management:} moisturiser, silicone products, and time
\end{itemize}

\section*{Complications}
\begin{itemize}
    \item Bacterial infection of the broken skin, including cellulitis
    \item Scarring and permanent pigment changes
    \item Persistent itching and sleep disturbance
    \item In skin-picking disorder: significant tissue damage, distress, and social avoidance
    \item Reduced quality of life
\end{itemize}

\section*{Prognosis and Outcomes}
Excoriation heals well when the underlying itch or picking is treated. Itchy skin conditions respond to appropriate skincare, and most skin damage heals within weeks, though scars may persist. Skin-picking disorder improves with psychological therapy, and many people learn to control the behaviour and heal their skin.

\section*{Prevention and Self-Care}
\begin{itemize}
    \item Keep the skin moisturised and treat itching promptly
    \item Keep nails short and consider gloves at night
    \item Use cooling and distraction when the urge to scratch is strong
    \item Cover healing skin to reduce temptation to scratch
    \item For skin picking: notice the triggers and substitute a different behaviour
    \item Seek help for the underlying skin condition or mental health concern
\end{itemize}

\section*{Sources and Further Reading}
The following reputable sources provide further information for healthcare students and patients seeking reliable, current advice about excoriation and skin picking.
\begin{itemize}
    \item DermNet. \textit{Excoriation and skin picking.} https://dermnetnz.org/
    \item Healthdirect Australia. \textit{Skin conditions and itching.} https://www.healthdirect.gov.au/
    \item Eczema Association of Australasia. \textit{Itch and scratch management.} https://www.eczema.org.au/
    \item OCD and Anxiety support services. \textit{Excoriation disorder resources.} https://iocdf.org/
    \item NHS. \textit{Excoriation disorder.} https://www.nhs.uk/mental-health/
    \item MSD Manual. \textit{Excoriation disorder.} https://www.msdmanuals.com/
\end{itemize}
